\section{Conclusion (Заключение)}
\label{sec:conclusion}

В данной работе исследована фундаментальная проблема рассогласования классической задачи Maximum Likelihood Estimation (MLE) и удовлетворённости пользователя в музыкальных стриминговых сервисах (\textit{The Skip Dilemma}). 

Основные научно-практические выводы исследования:
\begin{enumerate}
    \item \textbf{Преодоление налога на выравнивание (Zero Alignment Tax):} Прямая оптимизация предпочтений (DPO) со вспомогательной SFT-регуляризацией ($\gamma=0.10$) доказала, что выравнивание модели под длительность прослушивания не требует жертвовать точностью ранжирования следующего трека. На отложенном тестовом множестве из 9\,048 пользователей датасета Yambda-50M достигнут статистически достоверный прирост \textbf{+0.67\% HR@10} ($p < 0.0001$, чистый выигрыш $+60$ пользователей по тесту МакНемара).
    \item \textbf{Превосходство DPO над двухстадийными методами:} Сквозное выравнивание всех 43.2M параметров Causal GPT-2 превзошло как 2-стадийное реранжирование на базе критика $Q(s, a)$, так и Actor-Critic PPO, устранив эффект Reward Hacking и проблему информационного плато замороженных эмбеддингов.
    \item \textbf{Оздоровление каталога:} Вопреки опасениям схлопывания политики в безопасные хиты, DPO снизило ранговую корреляцию Спирмена с популярностью на \textbf{$-12.1\%$} и расширило каталог активных рекомендаций на \textbf{+1\,027 уникальных треков} из длинного хвоста при росте энтропии Шеннона ($13.728 \to 13.815$ бит).
    \item \textbf{Авторегрессионная надёжность:} Мультишаговые роллауты подтвердили, что DPO не подвержено деградации от накопления ошибки (Exposure Bias) и обеспечивает \textbf{+1.5 секунды} прироста суммарного времени сессии на каждые 5 треков.
\end{enumerate}

\subsection{Инженерные рекомендации для продакшена}
Для команд рекомендательных систем индустриальных медиасервисов рекомендуются:
\begin{itemize}
    \item Использовать бинарную или пороговую стратегию майнинга предпочтений ($\Delta \ge 0.40$), устраняющую шум пограничных прослушиваний.
    \item Обязательно фиксировать SFT-якорь ($\gamma \in [0.05, 0.10]$) для удержания модели от коллапса языковой структуры.
    \item Отказаться от тяжелых двухстадийных оффлайн-критиков на замороженных эмбеддингах в пользу прямого DPO-согласования.
\end{itemize}

\subsection{Направления дальнейших исследований}
Перспективным вектором развития является интеграция мультимодальных аудио-эмбеддингов (MULE, MuLan, EnCodec) для холодного старта новых композиций, а также развертывание онлайн-бандитов с контекстным DPO в реальном времени.
